% Part: normal-modal-logic
% Chapter: completeness
% Section: frame-completeness

\documentclass[../../../include/open-logic-section]{subfiles}

\begin{document}

\olfileid{nml}{com}{fra}

\olsection{ફ્રેમ-પૂર્ણતા}

જ્યારે આપણે બતાવીએ કે આપેલા તર્કનું કૅનોનિકલ નિદર્શન અનુરૂપ ફ્રેમ
ગુણધર્મ ધરાવે છે, ત્યારે \Log{K}નો પૂર્ણતા પ્રમેય બીજા મોડલ તંત્રો
સુધી વિસ્તારી શકાય છે.

\begin{thm}\ollabel{thm:completeframeprops}
  જો સામાન્ય મોડલ તર્ક $\Sigma$માં \olref{tab:correspondencetable}ની
  ડાબી બાજુનાં !!{formula}s પૈકી કોઈ એક હોય, તો~$\Sigma$નું કૅનોનિકલ
  નિદર્શન જમણી બાજુનો અનુરૂપ ગુણધર્મ ધરાવે છે.
\end{thm}

\begin{table}[htp]
  \centering
    \begin{tabular}{| l || l |}
      \hline
      \emph{જો $\Sigma$માં \dots{} હોય} & \emph{તો $\Sigma$નું કૅનોનિકલ
        નિદર્શન આ પ્રકારનું છે:} \\
      \hline \hline
      \Ax{D}:\;\; $\Box!A \lif \Diamond !A$ & \quad સિરિયલ; \\
      \hline
      \Ax{T}:\;\; $\Box!A \lif !A$ &\quad સ્વપ્રતિબિંબિત;\\
      \hline
      \Ax{B}: \;\;$!A \lif \Box\Diamond!A$ &\quad સંમિત; \\
      \hline
      \Ax{4}: \;\; $\Box!A \lif \Box\Box!A$ & \quad પરંપરિત; \\
      \hline
      \Ax{5}: \;\; $\Diamond !A \lif \Box\Diamond!A$& \quad યુક્લિડિયન.\\
      \hline
    \end{tabular}
    \caption{મૂળભૂત અનુરૂપતાની હકીકતો.}
    \ollabel{tab:correspondencetable}
\end{table}

\begin{proof}
  આ દરેક કિસ્સો ક્રમે લઈએ.

  ધારો કે $\Sigma$માં \Ax{D} છે, અને $\Delta \in W^\Sigma$ લો.
  એવું $\Delta'$ છે કે $R^\Sigma \Delta\Delta'$ તે બતાવવું છે.
  $\Box^{-1}\Delta$ એ $\Sigma$-સુસંગત છે તે બતાવવું પૂરતું છે:
  ત્યાર પછી લિન્ડનબાઉમના સહાયક પ્રમેયથી એવો પૂર્ણ
  $\Sigma$-સુસંગત ગણ $\Delta' \supseteq \Box^{-1}\Delta$ મળે છે,
  અને $R^\Sigma$ની વ્યાખ્યા
  \iftag{prvBox}{}{અને \olref[mod]{lem:box-iff-diamond} }મુજબ
  $R^\Sigma \Delta\Delta'$. હવે વિરોધ માટે ધારો કે
  $\Box^{-1}\Delta$ એ $\Sigma$-સુસંગત \emph{નથી}, એટલે કે
  $\Box^{-1}\Delta \Proves[\Sigma] \lfalse$.
  \olref[mod]{lem:box2} મુજબ $\Delta \Proves[\Sigma]
  \Box\lfalse$; અને $\Sigma$માં \Ax{D} હોવાથી
  $\Delta \Proves[\Sigma] \Diamond\lfalse$ પણ મળે છે.
  પરંતુ $\Sigma$ સામાન્ય હોવાથી $\Sigma \Proves
  \lnot\Diamond \lfalse$ (\olref[prf][nor]{prop:notDiamondBot});
  તેથી $\Delta \Proves[\Sigma] \lnot\Diamond \lfalse$ પણ,
  જે~$\Delta$ની સુસંગતતાનો વિરોધ કરે છે.

  હવે ધારો કે $\Sigma$માં \Ax{T} છે, અને $\Delta \in W^\Sigma$
  લો. $R^\Sigma \Delta\Delta$, એટલે કે
  \iftag{prvBox}{}{\olref[mod]{lem:box-iff-diamond} મુજબ}%
  $\Box^{-1}\Delta \subseteq \Delta$, બતાવવું છે. પરંતુ
  $\Box!A \in \Delta$ હોય તો \Ax{T}થી $!A \in \Delta$ પણ,
  જેમ જોઈતું હતું.

  હવે ધારો કે $\Sigma$માં \Ax{B} છે, અને $\Delta$,
  $\Delta' \in W^\Sigma$ માટે $R^\Sigma \Delta\Delta'$.
  $R^\Sigma \Delta'\Delta$ બતાવવું છે, એટલે કે
  \iftag{prvBox}{$\Box^{-1}\Delta' \subseteq \Delta$.
    \olref[mod]{lem:box-iff-diamond} મુજબ આ સમકક્ષ છે}{}
  $\Diamond\Delta \subseteq \Delta'$. તેથી $!A \in \Delta$
  ધારો. \Ax{B}થી $\Box\Diamond!A \in \Delta$ પણ.
  $R^\Sigma \Delta\Delta'$ની પૂર્વધારણા
  \iftag{prvBox}{}{અને \olref[mod]{lem:box-iff-diamond}} મુજબ
  $\Box^{-1}\Delta \subseteq \Delta'$; તેથી જરૂરી મુજબ
  $\Diamond!A \in \Delta'$.

  હવે ધારો કે $\Sigma$માં \Ax{4} છે, અને $R^\Sigma
  \Delta_1\Delta_2$ તથા $R^\Sigma \Delta_2\Delta_3$.
  $R^\Sigma \Delta_1\Delta_3$ બતાવવું છે. પૂર્વધારણા
  \iftag{prvBox}{}{અને \olref[mod]{lem:box-iff-diamond}} મુજબ
  $\Box^{-1}\Delta_1 \subseteq \Delta_2$ અને
  $\Box^{-1}\Delta_2 \subseteq \Delta_3$ બંને મળે છે.
  $R^\Sigma \Delta_1\Delta_3$ માટે
  $\Box^{-1}\Delta_1 \subseteq \Delta_3$ બતાવવું પૂરતું છે.
  તેથી $!B \in \Box^{-1}\Delta_1$, એટલે કે
  $\Box!B \in \Delta_1$, લો. \Ax{4}થી
  $\Box\Box!B \in \Delta_1$ પણ; પછી પૂર્વધારણાથી પહેલાં
  $\Box!B \in \Delta_2$ અને ત્યાર બાદ $!B \in \Delta_3$,
  જેમ જોઈતું હતું.

  હવે ધારો કે $\Sigma$માં \Ax{5} છે, અને $R^\Sigma
  \Delta_1\Delta_2$ તથા $R^\Sigma \Delta_1\Delta_3$.
  $R^\Sigma \Delta_2\Delta_3$ બતાવવું છે. પહેલી પૂર્વધારણાથી
  $\Box^{-1}\Delta_1 \subseteq \Delta_2$
  \iftag{prvBox}{}{\olref[mod]{lem:box-iff-diamond} મુજબ} મળે છે;
  બીજી પૂર્વધારણા $\Diamond\Delta_3 \subseteq \Delta_1$ને
  સમકક્ષ છે\iftag{prvBox}{,
    \olref[mod]{lem:box-iff-diamond} મુજબ}{}.
  $R^\Sigma \Delta_2\Delta_3$ બતાવવા
  \iftag{prvBox}{\olref[mod]{lem:box-iff-diamond} મુજબ
    એટલું પૂરતું છે કે}{આપણે બતાવવું પડશે કે}%
  $\Diamond\Delta_3 \subseteq \Delta_2$.
  તેથી $\Diamond!A \in \Diamond\Delta_3$, એટલે કે
  $!A \in \Delta_3$, લો. બીજી પૂર્વધારણાથી
  $\Diamond!A \in \Delta_1$ અને \Ax{5}થી
  $\Box\Diamond!A \in \Delta_1$ પણ. હવે પહેલી પૂર્વધારણાથી
  $\Diamond!A \in \Delta_2$, જેમ જોઈતું હતું.
  \end{proof}

પરિણામ તરીકે અનેક તંત્રોના પૂર્ણતા-પરિણામો મળે છે. દાખલા તરીકે,
$\Log{S5} = \Log{KT5} = \Log{KTB4}$ એ બધાં સ્વપ્રતિબિંબિત અને
યુક્લિડિયન નિદર્શનોના વર્ગને અનુલક્ષીને પૂર્ણ છે; આ વર્ગ બધાં
સ્વપ્રતિબિંબિત, સંમિત અને પરંપરિત નિદર્શનોના વર્ગ સમાન છે.

\begin{thm}\ollabel{thm:generaldet}
  $\mClass{C}_\Ax{D}$, $\mClass{C}_\Ax{T}$,
  $\mClass{C}_\Ax{B}$, $\mClass{C}_\Ax{4}$ અને
  $\mClass{C}_\Ax{5}$ અનુક્રમે બધાં સિરિયલ, સ્વપ્રતિબિંબિત,
  સંમિત, પરંપરિત અને યુક્લિડિયન નિદર્શનોના વર્ગો હોય. ત્યારે
  \Ax{D}, \Ax{T}, \Ax{B}, \Ax{4} અને \Ax{5}માંથી કોઈપણ
  સ્કીમા $!A_1$, \dots, $!A_n$ માટે તંત્ર
  $\Log{K}!A_1 \dots !A_n$ નિદર્શનોના વર્ગ
  $\mClass{C} = \mClass{C}_{!A_1} \cap \dots
  \cap \mClass{C}_{!A_n}$ વડે નિર્ધારિત થાય છે.
\end{thm}

\begin{prop}
  ધારો કે $\Sigma$ સામાન્ય મોડલ તર્ક છે. ત્યારે:
  \begin{enumerate}
  \item\ollabel{prop:anotherfive-a}%
    જો $\Sigma$માં સ્કીમા $\Diamond!A \lif \Box
    !A$ હોય, તો $\Sigma$નું કૅનોનિકલ નિદર્શન આંશિક વિધેયાત્મક છે.
  \item જો $\Sigma$માં સ્કીમા $\Diamond!A \liff \Box
    !A$ હોય, તો $\Sigma$નું કૅનોનિકલ નિદર્શન વિધેયાત્મક છે.
  \item જો $\Sigma$માં સ્કીમા $\Box\Box!A \lif \Box
    !A$ હોય, તો $\Sigma$નું કૅનોનિકલ નિદર્શન નબળે સઘન છે.
  \end{enumerate}
(આ ફ્રેમ ગુણધર્મોની વ્યાખ્યાઓ માટે \olref[frd][acc]{tab:anotherfive}
  જુઓ).
\end{prop}

\begin{proof}
  \begin{enumerate}
  \item ધારો કે $\Sigma$માં સ્કીમા $\Diamond !A \lif
    \Box !A$ છે. $R^\Sigma$ આંશિક વિધેયાત્મક છે તે બતાવવા
    કોઈપણ $\Delta_1$, $\Delta_2$, $\Delta_3 \in W^\Sigma$
    માટે $R^\Sigma \Delta_1\Delta_2$ અને $R^\Sigma
    \Delta_1\Delta_3$ હોય તો $\Delta_2=\Delta_3$ તે સાબિત કરવું
    પડશે. $R^\Sigma \Delta_1\Delta_2$ હોવાથી
    $\Box^{-1}\Delta_1 \subseteq \Delta_2$, અને
    $R^\Sigma \Delta_1\Delta_3$ હોવાથી
    $\Box^{-1}\Delta_1 \subseteq \Delta_3$ પણ
    \iftag{prvBox}{}{, બંને \olref[mod]{lem:box-iff-diamond} મુજબ}.
    $\Delta_2=\Delta_3$ માટે બે સમાવેશ
    $\Delta_2 \subseteq \Delta_3$ અને $\Delta_3 \subseteq
    \Delta_2$ સ્થાપિત કરવાનું પૂરતું છે. પહેલા સમાવેશ માટે
    $!A \in \Delta_2$ લો; ત્યારે $\Diamond!A \in \Delta_1$,
    અને સ્કીમા તથા $\Delta_1$ની નિષ્પત્તિ હેઠળની બંધતાથી
    $\Box!A \in \Delta_1$ પણ. તેથી $R^\Sigma
    \Delta_1\Delta_3$ની પૂર્વધારણાથી $!A \in \Delta_3$.
    બીજો સમાવેશ એ જ રીતે મળે છે.

  \item આ મુદ્દા~\olref{prop:anotherfive-a} અને
    \olref{thm:completeframeprops}માં સિરિયલતાના પ્રમાણમાંથી
    તરત ફલિત થાય છે.

  \item ધારો કે $\Sigma$માં સ્કીમા $\Box\Box!A \lif \Box!A$
    છે. $R^\Sigma$ નબળે સઘન છે તે બતાવવા $R^\Sigma
    \Delta_1\Delta_2$ લો. એવું પૂર્ણ $\Sigma$-સુસંગત
    ગણ $\Delta_3$ છે કે $R^\Sigma \Delta_1\Delta_3$ અને
    $R^\Sigma \Delta_3\Delta_2$ તે બતાવવું છે. મૂકો:
    \[
    \Gamma = \Box^{-1}\Delta_1 \cup \Diamond\Delta_2.
    \]
    $\Gamma$ એ $\Sigma$-સુસંગત છે તે બતાવવું પૂરતું છે;
    ત્યારે લિન્ડનબાઉમના સહાયક પ્રમેયથી તેને એવા પૂર્ણ
    $\Sigma$-સુસંગત ગણ~$\Delta_3$ સુધી વિસ્તારી શકાય કે
    $\Box^{-1}\Delta_1 \subseteq \Delta_3$ અને
    $\Diamond\Delta_2 \subseteq \Delta_3$, એટલે કે
    $R^\Sigma \Delta_1\Delta_3$\iftag{prvBox}{}{ (
      \olref[mod]{lem:box-iff-diamond} મુજબ)} અને $R^\Sigma
    \Delta_3\Delta_2$\iftag{prvBox}{ (
      \olref[mod]{lem:box-iff-diamond} મુજબ)}{}.

    વિરોધ માટે ધારો કે $\Gamma$ સુસંગત નથી. ત્યારે એવાં
    !!{formula}s $\Box!A_1$, \dots, $\Box!A_n \in \Delta_1$
    અને $!B_1$, \dots,~$!B_m \in \Delta_2$ છે કે
    \[!A_1, \dots,
    !A_n, \Diamond!B_1, \dots, \Diamond!B_m \Proves[\Sigma]
    \lfalse.\]
    $\Diamond (!B_1 \land \dots \land !B_m) \to
    (\Diamond!B_1 \land \dots \land \Diamond!B_m)$ દરેક સામાન્ય
    મોડલ તર્કમાં !!{derivable} હોવાથી આપણે નીચે મુજબ દલીલ કરીને
    $\Delta_2$ની સુસંગતતાનો વિરોધ મેળવીએ છીએ:
    \begin{align*}
      !A_1, \dots, !A_n, & \Diamond!B_1, \dots,
      \Diamond!B_m \Proves[\Sigma] \lfalse \\
      !A_1, \dots,!A_n & \Proves[\Sigma] (\Diamond!B_1 \land
      \dots \land \Diamond!B_m) \lif \lfalse\\
      & \qquad\text{નિષ્પત્તિ પ્રમેયથી}\\
      & \qquad\text{\olref[prf][prp]{prop:derivabilityfacts}\olref[prf][prp]{prop:derivabilityfacts-deduction}, અને \Taut} \\
      !A_1, \dots,!A_n & \Proves[\Sigma]
      \Diamond(!B_1 \land
      \dots \land !B_m) \lif \lfalse\\
      & \qquad \text{કારણ કે $\Sigma$ સામાન્ય છે} \\
      !A_1, \dots,!A_n & \Proves[\Sigma]
      \lnot\Diamond (!B_1 \land \dots \land !B_m)\\
      & \qquad\text{\PL\ મુજબ} \\
      !A_1, \dots,!A_n & \Proves[\Sigma]
      \Box\lnot (!B_1 \land \dots \land !B_m)\\
      & \qquad\text{$\lnot\Diamond$ને બદલે $\Box\lnot$} \\
      \Box!A_1, \dots,\Box!A_n
      & \Proves[\Sigma] \Box\Box \lnot (!B_1 \land \dots \land !B_m)\\
      & \qquad\text{\olref[mod]{lem:box1} મુજબ} \\
      \Box!A_1, \dots,\Box!A_n
      & \Proves[\Sigma]
      \Box\lnot (!B_1 \land \dots \land !B_m)\\
      &\qquad\text{સ્કીમા $\Box\Box!A \lif \Box!A$થી} \\
      \Delta_1 & \Proves[\Sigma]
      \Box\lnot (!B_1 \land \dots \land !B_m)\\
      &\qquad\text{એકદિશવર્ધિતાથી, \olref[prf][prp]{prop:derivabilityfacts}\olref[prf][prp]{prop:derivabilityfacts-monotonicity}} \\
      & \Box\lnot (!B_1 \land \dots \land !B_m) \in \Delta_1\\
      &\qquad\text{નિષ્પત્તિ હેઠળની બંધતાથી}; \\
      & \lnot (!B_1 \land \dots \land !B_m) \in \Delta_2\\
      &\qquad \text{કારણ કે }
      R^\Sigma \Delta_1\Delta_2.
    \end{align*}
    પરંતુ પસંદ કરેલાં બધાં સૂત્રો બીજા ગણમાં પણ છે, જેથી તેની
    સુસંગતતાનો વિરોધ મળે છે.
  \end{enumerate}
\end{proof}

આ ઉદાહરણોથી એવું લાગી શકે કે મોડલ તર્કનું દરેક તંત્ર~$\Sigma$
એ અર્થમાં \emph{પૂર્ણ} છે કે $\Sigma$ના દરેક પ્રમેય માન્ય હોય તેવી
દરેક ફ્રેમમાં માન્ય હોય તેવું દરેક સૂત્ર તે સાબિત કરે છે.
દુર્ભાગ્યે, આ અર્થમાં પૂર્ણ ન હોય એવાં ઘણાં તંત્રો છે.

\end{document}
